\section{Antecedentes}\label{sec:antecedentes}
\justifying

\subsection{Trabajos relacionados}

\subsubsection{Google Forms en la evaluación diagnóstica como apoyo en las actividades docentes. Caso con estudiantes de la Licenciatura en Turismo \cite{leyva2018google}.}

Google Forms atiende la necesidad de recopilar, organizar y analizar información de manera eficiente. Permite crear encuestas y formularios sin conocimientos técnicos, con almacenamiento automático en Google Sheets e integración con otros servicios de Google. Además, permite utilizar plantillas prediseñadas o imágenes y logos propios.


\subsubsection{La autoevaluación docente de aula: un camino para mejorar la práctica educativa \cite{echeverria2011autoevaluacion}.}

El artículo aborda la temática de la autoevaluación docente como estrategia que permite recoger información relevante acerca del desempeño docente del aula. Se analiza la autoevaluación docente como una estrategia constructiva de mejora profesional. Se promueve la reflexión crítica sobre el desempeño en el aula, enfocándose en identificar fortalezas y áreas de oportunidad.

\subsubsection{Sistema para la configuración de formularios y captura de los resultados de la evaluación de Atributos de Egreso en un proceso de mejora continua \cite{Cuautle2022}.}

 La evaluación de los Atributos de Egreso requiere información precisa en cada una de las evaluaciones. No obstante, al ser un proceso manual, es propenso a errores. El sistema automatiza la captura de información en procesos de evaluación académica. Incluye módulos que permiten precargar datos, validar usuarios y estructurar formularios según la UEA y los atributos a evaluar.

\subsubsection{Sistema académico de gestión escolar trimestralizado \cite{callisaya2024sistema}.}

Mediante un análisis preliminar realizado en la Unidad Educativa Técnico Humanístico "José Luis Suárez Guzmán" ubicada en la ciudad de El Alto, se observó que la institución no dispone de un sistema de monitoreo y registro virtual. El sistema propuesto centraliza la gestión académica, permitiendo consultar y registrar información pedagógica desde cualquier dispositivo. Mejora la trazabilidad y disponibilidad de los registros escolares.

\subsubsection{Sistema de Gestión Académica a través del desarrollo 
de Modelo-Vista-Controlador \cite{Rojas2019}.}

La gestión administrativa suele manejar grandes volúmenes de datos, que sin el tratamiento correcto crean dificultades en las tareas. Este proyecto implementa módulos específicos para la gestión académica: planes de estudio, carga académica y disponibilidad docente. Utiliza una arquitectura Modelo-Vista-Controlador (MVC) para mejorar la organización del desarrollo.

\subsubsection{Sistema de Gestión Administrativa para la Escuela Secundaria No. 44 "Jorge Luis Borges" \cite{gonzalez2017sistema}.}

El incremento significativo de la matrícula y del personal de la escuela produce un aumento en la cantidad de la documentación que se administra en forma manual, lo que origina errores en el registro de los datos. El sistema atiende el incremento de la carga administrativa,automatiza registros escolares y centraliza la información institucional. Fue construido con metodología XP, optimizando procesos en un entorno de cambio constante.

Las similitudes y diferencias con el proyecto propuesto se presentan en la Tabla \ref{table:trabajosRelacionados}.

%Esta es la tabla en donde se presentan las similitudes y diferencias.
\begin{longtable}{p{0.10\textwidth} p{0.40\textwidth} p{0.40\textwidth}}
  \label{table:trabajosRelacionados}\\
\hline
\textbf{Ref.}&
\textbf{Similitudes} &
\textbf{Diferencias} \\
\hline
\endfirsthead

\hline
\textbf{Ref.}&
\textbf{Similitudes} &
\textbf{Diferencias} \\
\hline
\endhead
\hline
\caption{Comparación cualitativa de los trabajos relacionados con el proyecto propuesto.}
\endlastfoot
\cite{leyva2018google} &
    \begin{itemize}[topsep=0pt,itemsep=0pt,parsep=0pt,partopsep=0pt]
        \item Uso de formularios digitales para capturar información.
        \item No requiere conocimientos técnicos.
    \end{itemize} &
    \begin{itemize}[topsep=0pt,itemsep=0pt,parsep=0pt,partopsep=0pt]
        \item Permite trabajar individualmente o de forma colaborativa a distancia.
        \item No tiene un sistema de autenticación con roles específicos.
    \end{itemize} \\
    \midrule

\cite{echeverria2011autoevaluacion} &
    \begin{itemize}[topsep=0pt,itemsep=0pt,parsep=0pt,partopsep=0pt]
        \item Relevancia de la autoevaluación docente como herramienta de mejora.
        \item Permitir que el docente identifique sus fortalezas y áreas de mejora.
    \end{itemize} &
    \begin{itemize}[topsep=0pt,itemsep=0pt,parsep=0pt,partopsep=0pt]
        \item Aborda la autoevaluación desde una perspectiva pedagógica, no tecnológica.
        \item No plantea un sistema automatizado.
    \end{itemize} \\
\midrule
    
\cite{Cuautle2022} &
    \begin{itemize}[topsep=0pt,itemsep=0pt,parsep=0pt,partopsep=0pt]
        \item Captura estructurada de datos mediante formularios.
        \item Autocompletado de formularios a partir de la información en una base de datos.
    \end{itemize} &
    \begin{itemize}[topsep=0pt,itemsep=0pt,parsep=0pt,partopsep=0pt]
        \item Enfocado en evaluar atributos de egreso por UEA.
        \item Módulo para indicar los atributos de egreso a evaluar.
        \item Módulo para asignar a un profesor a la UEA a evaluar.
    \end{itemize} \\
\midrule

\cite{callisaya2024sistema} &
\begin{itemize}[topsep=0pt,itemsep=0pt,parsep=0pt,partopsep=0pt]
        \item Registro y visualización de información académica.
\end{itemize} &
\begin{itemize}[topsep=0pt,itemsep=0pt,parsep=0pt,partopsep=0pt]
        \item Orientado al seguimiento pedagógico de estudiantes, no a documentos docentes.
\end{itemize} \\
\midrule

\cite{Rojas2019} &
\begin{itemize}[topsep=0pt,itemsep=0pt,parsep=0pt,partopsep=0pt]
    \item Módulos especializados para distintos procesos académicos.
    \item Manejo de  información  de los planes de estudios.
\end{itemize} &
\begin{itemize}[topsep=0pt,itemsep=0pt,parsep=0pt,partopsep=0pt]
	\item Gestión enfocada en planes de estudio y disponibilidad docente, no en formularios.
\end{itemize} \\
\midrule

\cite{gonzalez2017sistema} &
\begin{itemize}[topsep=0pt,itemsep=0pt,parsep=0pt,partopsep=0pt]
	\item Centralización de datos. 
	\item Automatización de tareas administrativas.
\end{itemize} &
\begin{itemize}[topsep=0pt,itemsep=0pt,parsep=0pt,partopsep=0pt]
	\item Diseñado principalmente para el personal administrativo, no para profesores.
\end{itemize} \\
    \bottomrule
\end{longtable}

\newpage
